\section{OBJETIVOS}

\subsection{Objetivo Principal}

Desarrollar e implementar un sistema integral de inteligencia urbana y análisis de patrones delictivos basado en algoritmos de \textbf{Aprendizaje Automático No Supervisado}, orientado a la detección autónoma de concentraciones críticas de delitos (\textit{hotspots}), la caracterización semántica del \textit{modus operandi}, el modelado probabilístico continuo del riesgo espacial y el aislamiento de eventos criminales anómalos, validado sobre el dataset masivo de referencia internacional de Chicago ($1.192.845$ registros procesados en capa Bronze y $798.639$ registros limpios en capa Silver) y el corpus nacional de seguridad ciudadana de Bolivia ($137$ crónicas estructuradas), integrando sus resultados en una plataforma interactiva web y cartográfica WebGL en 3D.

\subsection{Objetivos Específicos}

\begin{itemize}
    \item \textbf{Ingeniería y depuración de datos espaciotemporales y textuales:} Construir un pipeline reproducible de limpieza, auditoría y transformación que incorpore la conversión de coordenadas geográficas a radianes para métrica esférica de \textit{Haversine}, codificación de periodicidad temporal mediante funciones trigonométricas circulares $(\sin, \cos)$ de horas y días, y normalización de texto para vectorización TF-IDF.
    \item \textbf{Clustering geoespacial basado en densidad y multiescala:} Implementar y comparar los algoritmos DBSCAN, HDBSCAN y OPTICS para descubrir agrupaciones de morfología arbitraria, gestionar densidades variables y aislar formalmente los incidentes delictivos esporádicos clasificados como ruido estadístico.
    \item \textbf{Minería semántica de lenguaje natural y arquetipos de crimen:} Integrar el algoritmo K-Means con representaciones matriciales TF-IDF ponderadas y reducción TruncatedSVD para descubrir los clústeres semánticos de \textit{modus operandi} delictivo y habilitar un clasificador en tiempo real.
    \item \textbf{Modelado de riesgo continuo mediante mixturas gaussianas (GMM):} Cuantificar la incertidumbre espaciotemporal del crimen ajustando modelos de densidad probabilística con matrices de covarianza completa mediante el algoritmo Esperanza-Maximización (EM), derivando elipses de riesgo a $2\sigma$ ($95.4\%$ de masa de probabilidad).
    \item \textbf{Taxonomía distrital mediante clustering jerárquico:} Estructurar dendrogramas aglomerativos con enlace de mínima varianza de Ward sobre matrices normalizadas de incidentes por distrito policial, validando su consistencia métrica mediante el coeficiente de correlación cofenética.
    \item \textbf{Detección no supervisada de anomalías con Autoencoders neuronales:} Diseñar y entrenar una arquitectura de red neuronal con cuello de botella ($D \to 16 \to 4 \to 16 \to D$) que compute el error cuadrático medio de reconstrucción (MSE) para individualizar delitos atípicos o casos criminales excepcionales.
    \item \textbf{Proyección en espacios latentes de dimensionalidad reducida:} Comparar técnicas de reducción de dimensionalidad lineales (PCA, TruncatedSVD, PPCA) y no lineales (t-SNE, UMAP) para proyectar espacios de alta dimensionalidad sobre variedades visualizables en dos y tres dimensiones.
    \item \textbf{Despliegue operativo y visualización inmersiva en Globo 3D:} Construir una arquitectura desacoplada basada en microservicios Flask (API REST) y un cliente interactivo WebGL con MapLibre GL / Deck.gl que permita explorar dinámicamente los modelos analíticos.
\end{itemize}

\subsection{Árbol de Objetivos}

El Árbol de Objetivos presentado en la Figura \ref{fig:arbol_objetivos} transforma de forma biyectiva las causas del árbol de problemas en medios positivos de ejecución, y las consecuencias negativas en fines estratégicos alcanzados:

\begin{figure}[htbp]
\centering
\begin{tikzpicture}[
    node distance=1.0cm and 0.35cm,
    box_objetivo/.style={rectangle, draw=accentgreen!80!black, fill=accentgreen!10, text width=10.2cm, align=center, font=\bfseries\scriptsize, minimum height=1.3cm, rounded corners=3pt, line width=1.2pt},
    box_fin/.style={rectangle, draw=primaryblue!80, fill=primaryblue!8, text width=3.8cm, align=center, font=\tiny, minimum height=1.35cm, rounded corners=2pt, line width=0.8pt},
    box_fin_sup/.style={rectangle, draw=darkgray, fill=lightgray, text width=4.3cm, align=center, font=\tiny, minimum height=1.35cm, rounded corners=2pt, line width=0.8pt},
    box_medio_1/.style={rectangle, draw=secondaryblue!90, fill=secondaryblue!10, text width=3.8cm, align=center, font=\tiny, minimum height=1.45cm, rounded corners=2pt, line width=0.8pt},
    box_medio_2/.style={rectangle, draw=darkgray!70, fill=white, text width=3.8cm, align=center, font=\tiny, minimum height=1.45cm, rounded corners=2pt, line width=0.6pt},
    arrow/.style={->, >=stealth, thick, color=darkgray!80}
]

    % Objetivo Central
    \node (objetivo) [box_objetivo] {OBJETIVO CENTRAL:\\Desarrollar un sistema de Machine Learning No Supervisado para la delimitación geodésica de micro-hotspots, análisis semántico de modus operandi y aislamiento de anomalías espaciotemporales.};

    % Fines Directos
    \node (fin2) [box_fin, above=0.9cm of objetivo] {Mitigar la disparidad del 35\% en la cobertura de patrullaje preventivo mediante asignación guiada por densidad real.};
    \node (fin1) [box_fin, left=0.45cm of fin2] {Reducir la latencia analítica a tiempos casi reales en la correlación interjurisdiccional de incidentes emergentes.};
    \node (fin3) [box_fin, right=0.45cm of fin2] {Valorizar y estructurar computacionalmente más del 80\% de variables cualitativas de modus operandi a partir de texto libre.};

    % Fines Superiores
    \node (finsup1) [box_fin_sup, above=0.75cm of fin1] {Identificación y cuantificación continua de gradientes de riesgo en la evaluación longitudinal de series temporales urbanas.};
    \node (finsup2) [box_fin_sup, above=0.75cm of fin3] {Superación de los sesgos de agregación territorial en sistemas de información mediante modelado continuo georreferenciado.};

    % Conexiones Fines
    \draw [arrow] (objetivo) -- (fin1);
    \draw [arrow] (objetivo) -- (fin2);
    \draw [arrow] (objetivo) -- (fin3);
    \draw [arrow] (fin1) -- (finsup1);
    \draw [arrow] (fin3) -- (finsup2);

    % Medios Directos
    \node (medio2) [box_medio_1, below=0.9cm of objetivo] {Implementación de modelos de agrupamiento por densidad (DBSCAN, HDBSCAN, OPTICS) adaptables a morfologías viales complejas.};
    \node (medio1) [box_medio_1, left=0.45cm of medio2] {Aplicación rigurosa de métricas geodésicas ortodrómicas (Haversine) sobre coordenadas esferoidales $(\text{Lat}, \text{Lon})$.};
    \node (medio3) [box_medio_1, right=0.45cm of medio2] {Extracción y modelado semántico de crónicas narrativas mediante vectorización léxica (TF-IDF) y reducción dimensional latente.};

    % Medios Subyacentes
    \node (subm1) [box_medio_2, below=0.75cm of medio1] {Pipeline automatizado de preprocesamiento, tipado eficiente y conversión a Snappy Parquet.};
    \node (subm2) [box_medio_2, below=0.75cm of medio2] {Estimación continua de superficies de riesgo con Gaussian Mixture Models y elipses $2\sigma$.};
    \node (subm3) [box_medio_2, below=0.75cm of medio3] {Arquitectura desacoplada en Flask API REST y visualizador cartográfico interactivo Globo 3D.};

    % Conexiones Medios
    \draw [arrow] (medio1) -- (objetivo);
    \draw [arrow] (medio2) -- (objetivo);
    \draw [arrow] (medio3) -- (objetivo);
    \draw [arrow] (subm1) -- (medio1);
    \draw [arrow] (subm2) -- (medio2);
    \draw [arrow] (subm3) -- (medio3);

\end{tikzpicture}
\caption{Árbol de Objetivos del Proyecto de Inteligencia Urbana y Criminología Computacional.}
\label{fig:arbol_objetivos}
\figurenote{Elaboración propia a partir de la transposición biyectiva de las causas en medios de ejecución y los efectos negativos en fines estratégicos alcanzados.}
\end{figure}

